\subsection{First feedback and complete-execution exposure}
Table~\ref{tab:phase-complete} places the phase-level intervention beside the complete-execution endpoints. The six-class majority reference is 20/120 locally. Under Full forwarding, first post-ticket LR scores are 43/120 for Qwen and 40/120 for Llama, versus 25/120 and 19/120 before feedback. Projected forwarding gives 23/120 and 22/120. The Full--Projected point differences are 16.7 and 15.0 percentage points; unadjusted training/test-session refit intervals are $[4.0,29.8]$ and $[3.1,32.8]$ points.

All eight primary local comparisons exceed .05 after exact Holm adjustment, including every post-ticket--prefix contrast. Full--Projected LR values are .0778 and .1172; the smallest adjusted value, .0508, compares Llama's Full post-ticket and prefix phases. The unadjusted refit intervals answer a different, conditional uncertainty question. Exploratory RF scores are 50/120 versus 30/120 for Qwen and 59/120 versus 19/120 for Llama. CNN scores are 20/120 versus 20/120 and 26/120 versus 20/120. Table~\ref{tab:phase-complete} reports all four families. Their achieved scores provide attack evidence, not a leakage upper bound; the strongest test score does not replace the prespecified primary attacks.

\begin{table*}[t]\centering\small
\caption{Condition-matched recognition: correct predictions out of $n$ (six-class reference $n/6$). Prefix is the initial Full model call; Post is its first call after the ticket. Complete spans all requests and responses. RF and CNN are exploratory (E); neither replaces LR/k-FP in the primary family. $p_{proj}$ is the exact Holm-adjusted Full--Projected Post test. Complete-view comparisons are descriptive. Supplement S1 reports all uncertainty and frozen-transfer results.}\label{tab:phase-complete}
\setlength{\tabcolsep}{3pt}
\begin{tabularx}{\textwidth}{@{}lYrrrrrrrrr@{}}\toprule
Executor & Attack & $n$ & Prefix & \multicolumn{2}{c}{Post M} & \multicolumn{2}{c}{Complete M} & \multicolumn{2}{c}{Complete T+M} & $p_{proj}$ \\ & & & & Full & Proj. & Full & Proj. & Full & Proj. & \\\midrule
Qwen3:8B & LR & 120 & 25 & 43 & 23 & 46 & 44 & 54 & 39 & 0.0778 \\
Qwen3:8B & kFP & 120 & 18 & 32 & 20 & 43 & 39 & 39 & 30 & 0.1989 \\
Qwen3:8B & RF (E) & 120 & 25 & 50 & 30 & 61 & 31 & 55 & 30 & --- \\
Qwen3:8B & CNN (E) & 120 & 20 & 20 & 20 & 22 & 22 & 20 & 22 & --- \\
Llama3.1:8B & LR & 120 & 19 & 40 & 22 & 62 & 33 & 53 & 29 & 0.1172 \\
Llama3.1:8B & kFP & 120 & 14 & 31 & 20 & 47 & 34 & 45 & 31 & 0.1989 \\
Llama3.1:8B & RF (E) & 120 & 24 & 59 & 19 & 73 & 34 & 55 & 30 & --- \\
Llama3.1:8B & CNN (E) & 120 & 20 & 26 & 20 & 22 & 15 & 19 & 18 & --- \\
Hosted alias & LR & 48 & 5 & 15 & 10 & 9 & 9 & 12 & 6 & 0.9214 \\
Hosted alias & kFP & 48 & 3 & 6 & 7 & 9 & 5 & 13 & 11 & 1.0000 \\
Hosted alias & RF (E) & 48 & 8 & 15 & 9 & 7 & 7 & 13 & 8 & --- \\
Hosted alias & CNN (E) & 48 & 8 & 11 & 5 & 10 & 8 & 9 & 11 & --- \\
\bottomrule\end{tabularx}
\end{table*}


Complete-execution outcomes depend on the model and attack. Model-only LR changes from 46 to 44 correct for Qwen and from 62 to 33 for Llama; RF changes from 61 to 31 and from 73 to 34. Forwarding minimization leaves task-relevant fields and model responses in the execution. Their separate contributions remain unmeasured. Each local condition has 518 held-out tool calls, with equal call counts in every task pair; call order or action differs in 35 Qwen and 26 Llama pairs. Thus the contrast bundles content, encoded length and execution order, while the observed Llama drop is not explained by fewer tool calls. Supplement S1 gives the path audit and saved prefix diagnostics, including Qwen LR's unadjusted permutation $p=.0595$ for 25/120.

\subsection{Concurrency, event detection and cost}
Concurrency changes task association as well as traffic. Complete-window attacks are weaker under both Dedicated and Shared conditions, whose window representations pool all visible connections during overlapping executions. Their serial-to-concurrent decline consequently combines connection policy with window mixing.

Table~\ref{tab:randomized-events} reports the prespecified IoU .5 endpoint beside .3/.7 sensitivity checks, using the same saved event predictions. Under Shared, LR/RF recover 35.0--43.3\% of instances at .5 with 4.3--18.6\% precision and 8.21--44.79 false events per test session, including empty sessions. At .3, all-active exceeds both learned Shared detectors in F1 for both models: .295 versus .216/.106 for Qwen, and .508 versus .333/.464 for Llama. At .5, the learned detectors exceed that constant baseline. This combination of coarse predictions and a permissive criterion rewards session occupancy. At .5, tighter intervals matter, yet false alarms remain high. These are diagnostics of the selected detectors and overlap rule; they provide neither a deployment prior nor precision at a common recall target.

\begin{table*}[t]\centering\small
\caption{Event overlap sensitivity on the same predictions. Each condition contains 120 tasks in 24 sessions, including four empty sessions. Learned detectors use T+M; all-active emits one session interval per class. IoU .5 is primary; .3/.7 are sensitivity checks. FP/session divides all false events by 24, retaining empty sessions. Supplement S1 retains hourly normalization, counts and intervals for these bounded recordings.}\label{tab:randomized-events}
\setlength{\tabcolsep}{3pt}
\begin{tabularx}{\textwidth}{@{}llYrrrrrr@{}}\toprule
Executor & Policy & Detector & F1 (.3) & F1 (.5) & F1 (.7) & P (.5) & R (.5) & FP/session (.5) \\\midrule
Qwen3:8B & Shared & LR & 0.216 & 0.119 & 0.048 & 0.072 & 0.350 & 22.54 \\
Qwen3:8B & Shared & RF & 0.106 & 0.077 & 0.035 & 0.043 & 0.400 & 44.79 \\
Qwen3:8B & Shared & all-active & 0.295 & 0.038 & 0.000 & 0.035 & 0.042 & 5.79 \\
Qwen3:8B & Joint & LR & 0.196 & 0.117 & 0.050 & 0.070 & 0.375 & 25.04 \\
Qwen3:8B & Joint & RF & 0.329 & 0.188 & 0.090 & 0.123 & 0.400 & 14.29 \\
Qwen3:8B & Joint & all-active & 0.265 & 0.038 & 0.000 & 0.035 & 0.042 & 5.79 \\
Llama3.1:8B & Shared & LR & 0.333 & 0.216 & 0.104 & 0.144 & 0.433 & 12.88 \\
Llama3.1:8B & Shared & RF & 0.464 & 0.249 & 0.105 & 0.186 & 0.375 & 8.21 \\
Llama3.1:8B & Shared & all-active & 0.508 & 0.098 & 0.000 & 0.090 & 0.108 & 5.46 \\
Llama3.1:8B & Joint & LR & 0.441 & 0.251 & 0.104 & 0.186 & 0.383 & 8.38 \\
Llama3.1:8B & Joint & RF & 0.429 & 0.274 & 0.129 & 0.199 & 0.442 & 8.92 \\
Llama3.1:8B & Joint & all-active & 0.530 & 0.121 & 0.000 & 0.111 & 0.133 & 5.33 \\
\bottomrule\end{tabularx}
\end{table*}


Held-out goal completion is 119/120 for Qwen in every condition, and 117/120 for Llama except Projected at 118/120. Projection's model-byte ratio is about .78 of Full in both local cohorts. Joint shaping costs 19.5 and 22.5 times Shared's model-link bytes, while model-only LR changes from 25 to 18 and from 19 to 25 correct. Of 2,736 shaped model responses, 1,818 exceed the forty-frame minimum. Joint also improves Qwen's event RF score despite shaping (Table~\ref{tab:randomized-events}); the present contrasts do not isolate why. These mixed outcomes and variable lengths make this an observed design point, not a privacy--cost frontier. Supplement S1 gives complete completion, timing, bytes and uncertainty results.
